\chapter{文档理解与检索优化：PDF、表格、召回增强与 GraphRAG}

\begin{introduction}
\item {\heiti 本章核心问题}：当最小 RAG 已经跑通但效果仍不稳定时，应该优先优化文档理解、召回增强还是知识结构，判断顺序是什么？
\item {\heiti 主案例推进到哪一步}：把企业知识助手从“能检索”推进到“能稳定检索”，开始处理真实企业文档脏、乱、杂和口语化查询的问题。
\item {\heiti 读完后能做什么}：看到召回片段不准、PDF 混乱、表格断裂或关系问题频发时，能先判断该修解析、切分、查询增强、重排序还是 GraphRAG。
\end{introduction}

\section{学习目标}
\begin{itemize}
  \item 理解 RAG 进入真实企业文档后最常见的输入侧与召回侧瓶颈。
  \item 掌握 PDF 解析、表格保留、查询增强、重排序、负样本和 GraphRAG 的职责边界。
  \item 能按“先修输入，再修召回，最后才加结构”的顺序设计优化路线，而不是盲目堆技术。
  \item 能把企业知识助手中的典型坏例子映射到明确的修复入口。
\end{itemize}

\section{问题背景}
上一章建立的是最小可用闭环，它默认文档还算干净、查询还算标准、片段也大致可读。但只要系统进入真实企业环境，情况会马上复杂起来：制度文件是双栏 PDF，操作手册里混有截图和表格，审批说明散落在正文和附件里，员工提问却是“这个报销到底还能不能走特批”这样高度口语化的句子。

这时团队最容易出现的误判，是把所有效果问题都归结为“embedding 不够强”或者“模型不够聪明”。实际上，很多 RAG 项目的瓶颈根本不在模型，而在输入被解析错、表格被切碎、召回候选顺序不对、或者用户语言和文档语言根本没对齐。第 8 章的目标，就是把这些优化入口重新分层，建立一条有顺序的优化栈。

\section{核心概念}
\subsection{优化顺序比优化方法更重要}
RAG 优化最常见的失控方式，是团队听过什么方法就往系统里塞什么方法：今天上 HyDE，明天上重排序，后天又讨论 GraphRAG。这样做的问题不是方法本身没价值，而是大家往往不知道每种方法到底在修哪一层问题。

一个更稳定的优化顺序通常是：
\begin{enumerate}
  \item 先修文档理解，确保输入文本本身可读、可定位；
  \item 再修切分与结构保留，确保片段仍然是完整语义单元；
  \item 再修查询表达与召回排序，确保用户问题能找到正确证据；
  \item 只有在问题天然依赖实体关系和多跳链路时，才考虑更重的结构化方案。
\end{enumerate}

这条顺序的底层逻辑很简单：如果输入已经坏了，后面所有召回优化都是在坏文本上调参；如果候选已经对了，只是排不准，重排序比换一整套图结构更划算；如果问题本质只是“文档里怎么写”，那就没必要过早把系统推向 GraphRAG 复杂度。

\subsection{文档理解先于检索优化}
对很多企业场景来说，真正的第一道瓶颈不是检索器，而是\textbf{文档能否被正确读出来}。尤其是 PDF，它本质上更像一种排版容器，而不是天然的语义结构。单栏和双栏顺序不同，页眉页脚会混进正文，表格和注释可能被 OCR 打散，截图中的关键信息甚至根本不会进入文本层。

因此，只要召回片段本身已经是脏的、断裂的、带页眉页脚的，团队就应该先停下来修文档理解，而不是继续调 \texttt{top-k} 或换 embedding。修输入常常比修模型更能立刻提升系统上限。

\subsection{为什么 PDF 与表格会单独成为问题}
PDF 和表格之所以值得单独讲，不是因为它们“格式特殊”，而是因为它们携带的语义结构不同于普通段落文本。

PDF 的问题主要在\textbf{阅读顺序和版面结构}。同一页上的文字并不天然按阅读顺序排列，双栏、脚注、图注和侧边说明很容易打乱上下文。表格的问题主要在\textbf{行列关系}。当系统只按字符或句子切分时，表头、列名、数值和限制条件之间的绑定关系会被切断，结果就是检索到了“看起来相关的字”，却没有检索到真正有意义的结构。

这意味着：只要问题答案依赖“表头属于哪一列”“限制条件挂在哪一行”“附件中的金额对应哪个审批级别”，普通段落式切分就很可能不够用。

\subsection{标题树是最便宜也最常被忽视的结构恢复}
很多企业文档里，最先丢掉、却最值得保住的，不是正文句子，而是标题层级。因为对制度、手册和 FAQ 来说，标题树本身就在表达“这段规则属于哪个主题、适用于哪个范围、和前后哪个条款是一组”。如果解析阶段把这一层抹平，后面的切分和检索都会变得更盲。

标题树的重要性主要体现在三个方面。第一，它决定片段的\textbf{语义父节点}。同一句“需审批后执行”，挂在“差旅超标例外”下面和挂在“办公用品领用流程”下面，含义完全不同。第二，它决定片段的\textbf{引用路径}。真正上线时，用户更愿意看到“制度名称 > 第三章 > 住宿标准 > 例外审批”这类路径，而不是一段孤立正文。第三，它天然构成了一个很轻量的\textbf{结构索引}，很多问题其实先命中标题层，再回到正文块就够了。

所以，对大量制度型文档来说，标题提取往往是比“换一个更强 embedding”更划算的第一步。因为标题一旦保住，切分边界、二级索引、引用路径和后续摘要生成都会更稳。

\subsection{表格检索不要伪装成普通段落检索}
很多 RAG 系统在表格问答上长期效果差，并不是因为模型不会算，而是因为系统把表格先拍扁成了一段普通文本，再指望检索器自己还原行列关系。对工程团队来说，更稳的做法通常是把表格看成一条\textbf{单独的检索路径}。

这条路径至少要保住四类结构：
\begin{itemize}
  \item \textbf{表头}：告诉模型每一列到底代表什么。
  \item \textbf{行主键}：例如岗位、审批等级、城市档位、产品型号。
  \item \textbf{数值与单位}：金额、比例、时间长度、版本号不能脱离单位独立存在。
  \item \textbf{限制条件}：备注列、适用范围、例外条款和脚注常常比主体数值更关键。
\end{itemize}

\begin{table}[H]
\centering
\small
\begin{tabularx}{\textwidth}{>{\raggedright\arraybackslash}p{2.9cm}>{\raggedright\arraybackslash}p{3cm}>{\raggedright\arraybackslash}X}
\toprule
\textbf{做法} & \textbf{短期看起来的好处} & \textbf{长期真实代价} \\
\midrule
把整张表压成一段文本 & 实现简单，能快速进向量库 & 表头、单位、条件容易断裂，召回到也难用 \\
按行块保留表头绑定 & 结构更稳定，便于引用具体行列 & 需要单独设计切分与展示逻辑 \\
表格走独立索引或路由 & 更容易先命中正确结构对象 & 系统复杂度上升，需要问题分类与双路评测 \\
\bottomrule
\end{tabularx}
\caption{表格检索的难点不是“表格里也有字”，而是它有普通段落没有的结构}
\label{tab:table-retrieval-path}
\end{table}

因此，企业知识助手如果经常回答报销标准、价格矩阵、权限清单、排班表或 SLA 等问题，最好在前置分类时就判断“这是不是表格型问题”。一旦判断为表格型，就优先检索表对象、行对象或带表头绑定的片段，而不是让段落检索和重排序去硬猜哪些数字属于同一行。

\subsection{解析路线怎么选：规则提取、OCR 与视觉模型各有边界}
旧稿在 PDF 与表格解析上其实给过很清楚的路线差异，这里需要把它教材化。对工程团队来说，第一步不是问“最先进的模型是什么”，而是问\textbf{这份文档到底属于哪一种可读性条件。}

\begin{table}[H]
\centering
\small
\begin{tabularx}{\textwidth}{>{\raggedright\arraybackslash}p{2.8cm}>{\raggedright\arraybackslash}p{3.1cm}>{\raggedright\arraybackslash}X}
\toprule
\textbf{解析路线} & \textbf{更适合什么输入} & \textbf{主要提醒} \\
\midrule
规则型文本提取 & 原生电子 PDF、结构较规整的制度文档 & 速度快，但一旦阅读顺序复杂、标题层级混乱，脏片段会被成批放大 \\
OCR + 版面分析 & 扫描件、拍照件、缺少文本层的材料 & 先解决“看不见字”，但要警惕错字、漏字和版面重排错误 \\
视觉文档模型 & 复杂表格、图文混排、强依赖版面关系的页面 & 结构恢复更强，但成本更高，也更需要样本与评测支撑 \\
\bottomrule
\end{tabularx}
\caption{文档解析路线首先由输入条件决定，而不是由技术新旧决定}
\label{tab:document-parsing-routes}
\end{table}

这张表背后的工程结论很朴素：\textbf{先让文档被正确读出，再谈如何把它检索得更准。} 如果输入本来就是高质量原生 PDF，过早上重视觉模型未必划算；反过来，如果全是扫描件和复杂表格，死守纯规则抽取只会让后面每一轮检索优化都建立在坏文本上。

\subsection{双栏 PDF 为什么会把证据顺序读反}
双栏 PDF 是企业知识库里特别容易被低估的一类坑。对人眼来说，双栏页面的阅读顺序很自然：先读左栏，再读右栏；但对很多解析工具来说，页面上的文本框只是若干坐标对象。只要工具没有正确恢复阅读顺序，系统就可能把“左栏第一段 + 右栏第一段 + 左栏第二段”串成一条怪异文本。

这类错误的危险，不在于它会让片段看起来“完全乱码”，而在于它往往\textbf{半对半错}。也就是说，系统仍然能读出许多正确词语和句子，但关键信息的邻接关系已经被打乱。此时检索器会把这些片段判成“相关”，模型也会从中抽取出一些像样的关键词，最后却给出逻辑接错的答案。对企业问答来说，这种“看起来有依据”的错答尤其危险。

因此，遇到制度 PDF 时，很值得先做三类抽检：抽查若干页的左右栏顺序是否被打乱；抽查跨页标题是否还能挂回正确正文；抽查页眉页脚和脚注是否混入主段落。只要这三类问题明显存在，就该优先修阅读顺序，而不是先上 HyDE 或重排序。

\begin{table}[H]
\centering
\small
\begin{tabularx}{\textwidth}{>{\raggedright\arraybackslash}p{2.7cm}>{\raggedright\arraybackslash}p{3cm}>{\raggedright\arraybackslash}X}
\toprule
\textbf{文档类型} & \textbf{优先处理动作} & \textbf{为什么} \\
\midrule
原生电子 PDF & 保持标题层级与阅读顺序 & 通常文本可提取，关键是别把结构读乱 \\
扫描件 PDF & OCR 与版面识别优先 & 没有可靠文本层时，检索问题往往源于识别错误 \\
制度表格、报价表 & 保留表头与行列关系 & 数值和限制条件一旦失去绑定就无法准确作答 \\
FAQ 与网页知识页 & 保留问答块和标题路径 & 检索最依赖问题与答案边界，而不是字符长度平均 \\
\bottomrule
\end{tabularx}
\caption{不同知识载体进入 RAG 前的首要处理动作}
\label{tab:rag-doc-first-action}
\end{table}

\subsection{文档规范化、去重与版本淘汰：别让同一条制度以五种面貌入库}
很多团队把 RAG 优化理解成“如何更聪明地找”，却忽略了另一个同样关键的动作：\textbf{先让知识对象变得足够干净、足够唯一。} 企业知识库里，同一份制度往往会同时以 PDF 正式版、Word 编辑版、FAQ 摘要版、会议宣讲版和截图转发版多种形态存在。如果这些版本没有被规范化、去重和淘汰，后面再好的检索和重排序，也只能在一堆彼此近似却口径不一的候选里做更复杂的猜测。

更稳的做法通常包含三步。第一，给每份知识建立\textbf{canonical source}，也就是“真正权威的来源对象”；其他摘要、FAQ、截图或镜像版本都只作为辅助入口，而不是平等参与最终证据竞争。第二，建立\textbf{版本关系}，让系统知道哪些文档是替代关系、哪些是补充关系。第三，在入库前做\textbf{近重复检测}，避免同一段条款被多个轻度改写版本反复写进候选空间。

\begin{table}[H]
\centering
\small
\begin{tabularx}{\textwidth}{>{\raggedright\arraybackslash}p{3cm}>{\raggedright\arraybackslash}p{3cm}>{\raggedright\arraybackslash}X}
\toprule
\textbf{治理动作} & \textbf{主要在防什么} & \textbf{如果不做，后面最容易出现什么问题} \\
\midrule
建立权威主来源 & 防止摘要版、截图版与正式制度平权竞争 & 系统引用“像答案”的材料，却不是“应负责的材料” \\
版本淘汰与替代关系维护 & 防止新旧制度同时入模 & 召回结果表面很全，实际却自带冲突与时效风险 \\
近重复检测与合并 & 防止同一条规则占满候选位 & 上下文被重复片段挤爆，真正互补证据进不来 \\
\bottomrule
\end{tabularx}
\caption{RAG 优化不只是在找得更准，也是在让知识对象先变得可治理}
\label{tab:document-canonicalization}
\end{table}

\subsection{结构保留比“切得均匀”更重要}
在真实文档里，最有价值的语义单位往往不是“长度差不多的一段文本”，而是“一个规则块”“一张完整表”“一个 FAQ 问答对”“一组连续操作步骤”。因此，切分的目标从来不是让片段看起来整齐，而是让片段仍然像一个可独立理解的证据单元。

这也是为什么很多看似“更细致”的切分策略反而会伤害效果。片段变得更细，召回率表面上可能会上升，但真正关键的限定条件和例外条款被拆散后，模型接收到的是不完整证据，答案反而更容易偏。

企业知识助手尤其要警惕两种切坏方式。第一种是把条款结论和适用条件分开。第二种是把表头和数据行分开。前者会让模型过度泛化制度；后者会让模型误读金额、时限和审批层级。

\subsection{不同文档类型要用不同切分器，而不是一把尺子切到底}
旧稿在文本分块部分列过很多拆分方式，教材版里最值得保留的，不是具体库名，而是\textbf{文档类型决定切分器类型}这件事。FAQ、Markdown、HTML、代码片段、制度 PDF 和表格，本来就不是同一种语义对象。

一个够用的工程判断可以写成：
\begin{itemize}
  \item \textbf{FAQ / 问答页}：优先按问答对或标题块切，不要平均切成字符片段。
  \item \textbf{制度 PDF / 规章文档}：优先保留标题路径、条款层级和例外条件归属。
  \item \textbf{HTML / Markdown 文档}：优先利用已有标题树、列表和代码块边界。
  \item \textbf{表格和结构化附件}：优先按表头 + 行块保留，而不是转成大段纯文本。
\end{itemize}

这部分看起来像预处理细节，实际上直接决定后面检索到的是不是一个“能被引用的证据单元”。如果切分阶段就把标题、问答边界、表头关系和代码块上下文全部抹平，检索再准，也只能召回被压扁后的证据。

\subsection{chunk overlap 不是越大越稳：它是在保留局部依赖，不是在重复灌水}
很多团队发现证据边界容易被切坏后，第一反应是把 \texttt{chunk\_overlap} 一路开大。这个动作有时确实能缓解边界丢失，但它的职责其实很具体：\textbf{只是在块与块之间保住局部依赖，避免关键信息刚好落在切口上。} 它不是为了让同一份材料在向量库里重复出现很多次。

在企业文档里，真正需要 overlap 保护的，通常是三类局部关系：条款结论与适用条件挨得很近，操作步骤的前后文互相依赖，以及表头、注释与首行数据之间的绑定。如果系统主要错在这些地方，适度 overlap 往往有帮助；但如果问题是标题层级丢了、表格被拍扁了、双栏顺序错了，那么把 overlap 开得再大，也只是在重复扩散同一种坏片段。

过大的 overlap 还会带来三个很实际的副作用：
\begin{itemize}
  \item \textbf{重复召回变多}：同一段规则被多个近邻块反复命中，表面上像“召回更全”，实际上是在挤占别的候选位置。
  \item \textbf{上下文噪声上升}：重排序和压缩阶段会反复看到高度相似的片段，真正的新证据反而更难进入最终上下文。
  \item \textbf{更新成本变高}：制度一旦改版，需要重刷的近似块更多，版本漂移也更难控。
\end{itemize}

所以 overlap 的工程用法更像一把手术刀，而不是止疼药。稳妥做法通常是：先按文档类型决定切分器，再只对容易断开的局部关系加适度 overlap，并在验证集里专门检查“重复召回率”和“关键信息是否仍然跨块断裂”。只有这样，overlap 才是在修证据边界，而不是在放大库内冗余。

\subsection{父子块与多粒度切分：先让粗块找方向，再让细块给证据}
单一粒度切分最大的矛盾是：块切得太大，检索容易把无关内容一起带进来；块切得太小，证据又可能只剩半句规则、半行表格或者一截缺条件的步骤。对长制度、手册和多页说明文档来说，更稳的办法往往不是继续争论“512 还是 1024”，而是承认\textbf{召回和作答需要的粒度本来就不完全相同。}

父子块或多粒度切分就是在解决这个矛盾。粗块负责保留主题范围，例如“差旅报销总则”“超标审批例外”“VPN 访问开通流程”；细块负责给出可以直接引用的证据句、规则行或参数项。检索时，系统可以先用粗块判断方向，再回到对应父块下的细粒度证据做精取。这样做的好处是，召回阶段不必在过细片段里盲撞，生成阶段也不必拿着一整页大块文本去硬找一行答案。

对企业知识助手来说，这种分层切分尤其适合两类材料。第一类是\textbf{制度文档}，因为用户常常先问一个宽问题，如“住宿超标怎么处理”，系统需要先定位到“差旅 > 住宿 > 例外审批”这一段，再从细块里抓出金额阈值、审批角色和附件要求。第二类是\textbf{产品或运维手册}，因为用户会问“证书更新失败怎么办”这类故障问题，系统需要先锁定到相关章节，再把真正的排障步骤和限制条件抽出来。

这里要注意一个常见误区：多粒度不是简单地把同一篇文档切成两套不同长度的块，然后全部平铺进同一个库。更合理的做法是让不同粒度承担不同职责，并保留父子映射关系。否则系统虽然“看起来有多种粒度”，实际上仍然是在海量候选里无差别搜索，复杂度上来了，信息组织却没有真正变好。

\subsection{查询增强是在弥合“用户语言”和“文档语言”的落差}
很多企业员工不会用制度文档的正式说法来提问。员工会说“出差住宿超标还能不能报”，文档里却写“住宿补助标准”“超标审批”和“例外报销流程”；员工会说“代码库权限怎么开”，文档里却写“仓库访问申请”“审批路径”和“角色授权”。如果直接用原始问题去检索，很可能只命中部分相关证据。

这时才轮到查询增强出场。它的职责不是创造新知识，而是把用户语言变得更像文档语言。常见做法包括：
\begin{itemize}
  \item \textbf{问题改写}：把口语问题转成更适合检索的规范表达；
  \item \textbf{HyDE}：先生成一段假想答案，再用这段答案去检索；
  \item \textbf{多查询扩展}：从一个问题生成多个检索变体；
  \item \textbf{RAG-Fusion}：对多查询结果做融合排序，减少单一路径遗漏。
\end{itemize}

这些方法的共同点是：它们都在修“表达错位”，而不是修“输入脏文本”。如果原文都已经被解析坏了，查询再聪明也很难从错误材料里找出正确答案。

\begin{table}[H]
\centering
\small
\begin{tabularx}{\textwidth}{>{\raggedright\arraybackslash}p{2.5cm}>{\raggedright\arraybackslash}p{2.8cm}>{\raggedright\arraybackslash}p{2.5cm}>{\raggedright\arraybackslash}X}
\toprule
\textbf{方法} & \textbf{主要修什么} & \textbf{代价} & \textbf{适合何时用} \\
\midrule
问题改写 & 把口语问题转成文档语言 & 低 & 用户表达随意，但文档本身较干净 \\
HyDE & 放大潜在语义空间 & 中 & 问题短、抽象、缺少检索关键词时 \\
多查询扩展 & 覆盖多个表述角度 & 中到高 & 单个问题常对应多种叫法或多份制度时 \\
RAG-Fusion & 融合多查询召回结果 & 高 & 多查询有效，但结果顺序不稳定时 \\
\bottomrule
\end{tabularx}
\caption{查询增强方法的职责、成本与适用时机}
\label{tab:rag-query-enhancement}
\end{table}

\subsection{RAG-Fusion 不只是多查几次：关键在融合时别丢用户意图}
旧稿在 RAG-Fusion 上讲得比当前新版更细，这里需要把它补回来。RAG-Fusion 的价值，不是简单地把一个问题扩写成多个问题，而是\textbf{在多条查询路径都可能有用时，尽量保住原始用户意图，同时减少单一路径漏召回。}

它通常包含两步：先生成若干个语义相近但检索角度不同的查询，再把多路召回结果做融合排序。很多工程实现会使用 RRF（Reciprocal Rank Fusion）这一类简单而稳健的融合策略，因为它能在不过分依赖某一路绝对分数的前提下，把多条路径共同认可的候选往前推。

但这类方法也有明确代价。查询一多，召回成本、重排成本和意图漂移风险都会上升。特别是在企业制度问答里，如果扩写查询过度追求覆盖面，反而可能把“原问题最关心的限制条件”冲淡。因此，RAG-Fusion 真正难的地方不是“能不能生成多个问题”，而是\textbf{能不能在扩展表达的同时，守住原问题的约束中心。}

\subsection{RRF 为什么常被用作 RAG-Fusion 的默认融合器}
旧稿里把 RRF（Reciprocal Rank Fusion）单独展开过，这一层值得保留，因为它解释了为什么 RAG-Fusion 在工程上常常“够简单也够稳”。RRF 的直觉是：不要求某一路查询给出特别精确的绝对分数，而是看一个候选在多路结果里\textbf{是否都排得不差}。常见写法可以记成
\[
\mathrm{score}(d)=\sum_i \frac{1}{k+\mathrm{rank}_i(d)},
\]
其中 \(\mathrm{rank}_i(d)\) 表示候选 \(d\) 在第 \(i\) 路查询结果中的名次，\(k\) 是一个平滑常数。

这个公式背后的工程含义很清楚：如果一个片段在多路查询里都能稳定进入前列，它大概率比“只在某一路偶然冲到第一”的片段更可靠。RRF 因此特别适合企业知识场景，因为这里的查询改写往往并不完美，多路结果之间也很难直接比较原始分数。与其纠结“不同 embedding 分数能不能横向比较”，不如先用名次把共同认可的候选往前推。

当然，RRF 也不是无代价。它会天然偏爱“多路都不算差”的候选，因此如果某个问题其实特别依赖一条非常精确的关键词检索路，RRF 也可能把那条路的强信号稀释掉。所以更稳的做法通常是：先确认多查询确实能补召回，再把 RRF 当默认融合器，而不是把它当作任何场景都必赢的后处理。

\subsection{混合检索与元数据过滤：关键词、版本号和权限约束不要全交给向量}
到了真实企业环境以后，很多失败并不是“语义没对上”，而是问题里混着\textbf{精确字符串信号和业务约束信号}。例如员工会问“2025 版差旅制度里 P6 级别去东京的住宿上限是多少”“错误码 E4032 是什么权限问题”“华东区合同模板最新版本在哪里”。这类问题里，版本号、错误码、区域名、岗位级别、产品型号都不是普通语义近邻能稳定处理好的内容。

这就是为什么很多系统最后会回到混合检索。向量检索负责抓语义相关性，关键词检索负责抓精确词项，元数据过滤负责收紧业务边界。三者分工清楚时，系统会比“全靠 embedding 猜”稳得多。

\begin{table}[H]
\centering
\small
\begin{tabularx}{\textwidth}{>{\raggedright\arraybackslash}p{3.1cm}>{\raggedright\arraybackslash}p{3.1cm}>{\raggedright\arraybackslash}X}
\toprule
\textbf{问题里最强的信号} & \textbf{优先补哪类能力} & \textbf{原因} \\
\midrule
口语描述、同义表达、模糊问法 & 向量检索或查询增强 & 需要弥合用户语言与文档语言的落差 \\
错误码、版本号、产品名、缩写词 & 关键词检索或混合检索 & 精确词项常常比纯语义近邻更可靠 \\
区域、岗位、时间、生效版本、权限范围 & 元数据过滤 & 这些是业务约束，不该交给模型在后面“自己猜” \\
同一问题同时依赖语义和精确字段 & 混合检索 + 过滤 + 重排序 & 先拉到正确池子，再让排序器决定哪条最能支撑答案 \\
\bottomrule
\end{tabularx}
\caption{当查询里同时出现语义信号与业务约束时，混合检索往往比单一路径更稳}
\label{tab:hybrid-retrieval-metadata}
\end{table}

元数据过滤尤其容易被低估。它看起来不像模型技术那么“高级”，但对企业知识助手来说，经常是最便宜也最有效的降噪手段。因为很多问答错误，根本不是系统没找到相关材料，而是它把\textbf{旧版本、别的业务线、别的地区或用户无权限查看的材料}也一起带了进来。把这些边界前移到检索阶段，通常比让大模型在一堆候选里事后分辨更稳。

这也解释了一个常见的工程原则：\textbf{能用结构化约束解决的问题，不要都推迟到生成阶段。} 如果用户明确问的是“最新版本”“某地区”“某岗位”或“某系统错误码”，那就应该尽量在检索入口把候选池收紧，而不是让模型在生成时顺手猜哪一份更像正确答案。

\subsection{FLARE 与前瞻检索：答案生成过程中也可能需要再补证据}
旧稿里的 FLARE 提醒我们：检索不一定只能发生在回答之前。有些复杂问题在第一轮检索时还看不出缺口，真正开始组织答案时，系统才会发现“这里少一条限定条件”“这里缺一个例外流程”“这里需要再补一列表格字段”。

这时更稳的思路，不是默认把第一轮召回硬塞到底，而是允许系统在生成过程中做一次前瞻判断：当前证据是否足以支持下一步回答。如果不足，就再触发一轮更有针对性的检索。这类方法的工程意义不在于“让链路更炫”，而在于减少两种常见失败：
\begin{itemize}
  \item 第一轮召回看起来相关，但真正下结论时证据不够；
  \item 模型已经开始往一个方向写，却没有意识到自己缺最后那一条关键条款。
\end{itemize}

当然，FLARE 一类方法也意味着更高延迟和更复杂的状态管理，所以它更适合高价值复杂问答，而不是所有请求的默认主链。

\subsection{检索器对齐：让 query、doc 与生成目标说同一种话}
很多检索问题并不是“用户没问清楚”，也不是“文档太脏”，而是 query、document 和最终回答需求根本不在同一个表示空间里。用户问的是口语句子，文档写的是制度标题，模型最后却被要求输出结构化结论和引用。如果中间没有做对齐，检索器就很容易在三个目标之间左右摇摆。

所谓检索器对齐，重点不是把 embedding 调得“更神”，而是让检索表示更贴近真实任务。常见做法包括：把历史问答日志中的问题和最终证据片段配成训练样本；为制度文档额外生成 FAQ 问句或摘要标题，作为更适合检索的侧车索引；在检索前先做问题分类、语言识别、表格/正文路由，让不同类型的问题走不同入口。这样做的本质，是减少“用户怎么问”和“知识怎么写”之间的表达断层。

在企业知识助手里，小模型也经常在这一层发挥作用。它们未必负责最终生成，但很适合做意图分类、语言检测、是否需要表格检索、是否需要权限校验前置判断等轻任务。把这些前置判断交给小模型，可以让大模型和重检索链只处理真正需要它们的问题，从而把召回对齐做得更细，也把整体成本压得更稳。

\subsection{重排序与负样本在修“相关但不够好”的候选}
当召回列表里已经常常出现正确文档，但不是排在最前面，或者排在前面的片段“看起来相关却不足以回答问题”，这就不再是简单的检索不到，而是\textbf{候选排序质量}问题。此时，重排序器通常比盲目增大 \texttt{top-k} 更有效。

而当这种“相关但无用”的误召回长期存在时，就要考虑更深一层：检索器是否真的学会区分正例和难负例。难负例指的是那些与问题表面上非常像、词也很像，但其实不能支持答案的片段。对企业制度问答来说，这类难负例非常常见，例如“审批说明”和“备案说明”“住宿标准”和“交通补助标准”。

负样本训练的意义，就在于让检索器学会拒绝这些“貌似相关”的片段。它并不是最先上的优化项，但当系统规模扩大、通用 embedding 已经压不住领域歧义时，它会变得很有价值。

\subsection{把 embedding 训对：相似度、对比学习、假负例与 SimANS}
旧稿在这一块其实讲得比新版更细。RAG 优化不只是“多加一个重排序器”，还包括把稠密检索器本身训对。对稠密检索来说，查询和文档会先各自编码成向量，再通过点积或余弦相似度做匹配。于是，真正的核心问题就变成了：\textbf{什么样的样本会把“该靠近的向量拉近、该分开的向量推远”。}

最常见的训练方式是对比学习。直觉上，它会把问题与正确片段视为正例，把问题与不支持答案的片段视为负例。这个思路看起来简单，但企业知识库里很快会遇到两个现实麻烦：
\begin{itemize}
  \item \textbf{假负例很多}：同一个制度可能在多个版本、多个模板或 FAQ 摘要里重复出现。若只按“不是当前正例”就一律当负例，模型会被教成把本该相近的片段推远。
  \item \textbf{难负例更有价值}：真正拉开检索器能力差距的，往往不是完全无关的噪声，而是“看起来很像、但并不能回答当前问题”的片段。
\end{itemize}

\begin{table}[H]
\centering
\small
\begin{tabularx}{\textwidth}{>{\raggedright\arraybackslash}p{3cm}>{\raggedright\arraybackslash}p{3.2cm}>{\raggedright\arraybackslash}X}
\toprule
\textbf{训练信号} & \textbf{主要价值} & \textbf{主要风险或提醒} \\
\midrule
正例对 & 让问题靠近真正支持答案的证据 & 正例质量差时，整个检索器会学偏 \\
批内负样本 & 训练便宜、扩展方便 & 企业语料里容易误伤同义片段或重复版本 \\
人工或规则挖的难负例 & 最能提升“相关但不对”的区分力 & 构造成本更高，需要持续维护 \\
SimANS 一类动态挖负例 & 能从当前检索器最容易混淆的候选中持续采样 & 若评估体系不稳，容易把采样噪声当真实收益 \\
\bottomrule
\end{tabularx}
\caption{检索器训练的关键不只是有负样本，而是负样本够不够像真问题}
\label{tab:dense-retrieval-training}
\end{table}

所以，检索器微调并不是“样本越多越好”，而是“样本边界越清楚越好”。如果团队已经能稳定拿到问题-证据配对，且失败样本主要集中在“召回了看似相关但并不支持答案的片段”，这时再去做难负例治理、对比学习或 SimANS 一类动态采样，收益才会比较明确。反过来，如果连 PDF 解析和结构切分都还不稳，过早做 embedding 微调，往往只是在坏输入上学坏模式。

\subsection{先上重排序还是先训检索器：看正确证据出现在哪一层}
很多团队走到这里时会问：既然召回不稳，到底该先上 cross-encoder 重排序，还是直接去微调 embedding 检索器？一个实用判断是：\textbf{先看正确证据有没有稳定出现在候选集合里。} 如果正确片段经常已经在前二十或前三十里，只是排不进前几位，那么优先上重排序通常更划算；如果正确片段连候选都经常进不来，才更像检索器表示空间本身有问题。

\begin{table}[H]
\centering
\small
\begin{tabularx}{\textwidth}{>{\raggedright\arraybackslash}p{3.2cm}>{\raggedright\arraybackslash}p{2.8cm}>{\raggedright\arraybackslash}X}
\toprule
\textbf{现象} & \textbf{优先动作} & \textbf{原因} \\
\midrule
正确证据常在候选里，但排不进前几位 & 先上重排序 & 说明粗召回已有基础，问题更像排序质量不足 \\
正确证据经常完全不在候选里 & 先查检索器、切分或查询表示 & 排序器无米可炊，先得让正确证据有机会出现 \\
不同说法下召回波动很大 & 先查查询增强或对齐数据 & 这更像表达空间错位，而不只是排序问题 \\
领域歧义长期存在，且已有问题-证据样本 & 再考虑微调检索器 & 这时样本能直接用于修表示边界，收益更明确 \\
\bottomrule
\end{tabularx}
\caption{RAG 优化里，重排序和检索器微调解决的不是同一层问题}
\label{tab:rerank-vs-retriever}
\end{table}

这个顺序之所以重要，是因为重排序通常更快见效，也更容易局部验证；检索器微调则更像底层改造，一旦样本有偏，可能把整个候选空间一起带歪。教材里最该留下的，不是哪种方法“更高级”，而是知道自己当前到底缺粗召回、缺细排序，还是缺表达对齐。

\subsection{Cross-Encoder 重排序为什么常常是“最贵但最值”的那一步}
旧稿里很多经验最后都会落到同一个动作上：在稠密召回之后，再加一层 cross-encoder 或更强的精排器。它之所以值钱，是因为它看的不再只是 query 向量和 doc 向量的粗近似，而是\textbf{把问题和候选片段一起读一遍，再判断这段材料到底能不能支持当前问题。}

对企业知识助手来说，这一步尤其适合处理“看起来相关，但真正答题时并不够用”的候选。比如“住宿标准”和“住宿超标审批流程”都与问题相关，但只有后者真正回答了“超标后怎么办”。Embedding 检索常能把二者都找回来，重排序则更有机会把真正可作答的证据提到前面。

\begin{table}[H]
\centering
\small
\begin{tabularx}{\textwidth}{>{\raggedright\arraybackslash}p{2.8cm}>{\raggedright\arraybackslash}p{3cm}>{\raggedright\arraybackslash}X}
\toprule
\textbf{做法} & \textbf{更适合什么时候启用} & \textbf{最该控制什么预算} \\
\midrule
只做稠密召回 & 问题简单、知识库干净、候选歧义低 & 保持低延迟，不要无意义扩大候选数 \\
召回后加轻量重排序 & 候选常出现但前排有明显噪声 & 控制待重排候选数，避免精排吞掉整体吞吐 \\
召回后加强精排器 & 高风险问题、制度例外题、表格和正文混合理解题 & 把计算预算留给真正高价值问题，而不是所有请求全量上重排 \\
\bottomrule
\end{tabularx}
\caption{重排序的关键不是“要不要上”，而是“给哪些问题上、重排多少候选”}
\label{tab:rerank-budget}
\end{table}

这一步和前面的查询增强并不冲突。查询增强负责把候选找回来，重排序负责在候选已经出现后判断“谁最像真正可作答的证据”。很多系统之所以到了中期阶段开始明显变稳，往往不是因为又多会了一种召回技巧，而是终于在这里补上了“相关不等于可回答”的判断层。

\subsection{上下文压缩：召回之后还要再做一次“证据降噪”}
很多团队以为“召回正确片段”以后就万事大吉，于是把前十几个候选一股脑塞进 prompt。问题在于，召回正确不等于入模干净。片段一多，噪声、重复、近义条款和无关附录也会一起进入上下文，最终拖慢 prefill，稀释关键证据，甚至把模型带偏。

因此，检索后的上下文压缩通常是第二次筛证据。它不是简单地把字数砍短，而是把真正支撑答案的那部分材料留下来。一个常见顺序是：先保留候选中最相关的 3--5 个证据块；再在块内抽取关键句、关键行或关键字段；如果仍然过长，再做带来源指针的摘要压缩。对表格问答来说，压缩的对象往往不是整张表，而是“表头 + 命中的相关行 + 与该行绑定的限制条件”。

这里最容易犯的错，是把压缩做成“重新写一段差不多的说明”，却丢掉原始证据锚点。更稳的做法是：\textbf{压缩可以减少噪声，但不能切断回到原文的路径。} 换句话说，压缩后的材料仍应带着页码、条款号、表格行列或原文片段引用，这样后续生成、评测和复盘才不会失去依据链。

\subsection{长上下文重排：关键证据不要埋在中间}
即使召回和压缩都做得不错，把材料放进 prompt 的顺序也仍然会影响最终答案。长上下文场景下，一个经常被观察到的现象是：模型往往更容易利用开头和结尾的内容，而对中间区域的关键信息利用不足。于是，系统明明已经召回了正确证据，却因为把它埋在一串相似候选的中段，导致答案仍然抓不住重点。

这就是长上下文重排存在的意义。它不是在重新决定“哪些候选相关”，而是在回答另一个问题：\textbf{哪些证据应该最先被模型看见，哪些证据应该作为补充放在后面，哪些重复或次要材料根本不该进来。} 对制度问答来说，通常应把最能直接回答问题的条款块放前；对复杂综合题，则可以把总则或结论性片段放前，再把例外条件、附录和补充定义按支持关系往后排。

更重要的是，长上下文重排常常要和压缩一起看。若没有先去重，重排只是把一堆相似片段换个顺序；若没有保留来源锚点，重排后的候选虽然“看上去更顺”，却会让后续校验和引用更难回溯。所以这一步的核心并不是“把最相关的都往前堆”，而是让 prompt 里的证据呈现出一种\textbf{先主证据、再补证据、最后留回溯路径}的组织方式。

\subsection{GraphRAG 适合强关系和多跳问题，不是通用终点}
GraphRAG 很容易因为概念新而被过度使用。实际上，它真正有价值的前提，是问题本身天然依赖实体、关系和多跳链路。例如“某审批步骤依赖哪一级负责人”“某系统告警会影响哪些下游服务”“某产品组件与哪些许可证绑定”，这些都明显带有图结构特征。

但如果问题本质上仍然是“文档里到底怎么写”，例如“差旅超标如何审批”“私有化部署支持哪些模块”，那么先把文档理解、结构保留和召回质量做好，通常比直接上 GraphRAG 更划算。图结构会带来实体抽取、关系构建、图更新和多路径排序等额外成本，只有当问题确实需要这些能力时，这笔成本才值得付。

\subsection{实体命名规范先于建图：别让同一个对象裂成三份}
GraphRAG 一旦进入实现阶段，最容易被低估的，不是图数据库本身，而是\textbf{实体命名规范}。企业内部同一个对象往往有多个叫法：员工会说“代码库权限”，制度里写“仓库访问授权”，系统表单又叫“Git 访问申请”。如果这些别名在建图前没有被归一或挂到同一 canonical id 上，图就会天然碎裂。

这类碎裂会直接伤害后续子图检索。用户问一个口语实体名，系统只命中其中一部分节点；关系边明明都存在，却被分散在多个近义实体上；最后检出来的子图看起来不算错，却永远缺一两条关键边。于是团队会误以为“图推理还不够强”，其实更前面的问题是实体标准化没做好。

因此，GraphRAG 的一个很朴素但很关键的前提是：在实体入图前先做别名表、规范名映射和版本统一。对企业知识助手来说，这甚至可以从最轻量的方式开始，例如给“VPN”“远程接入”“安全接入网关”先建同义词归并表。只有这一层打稳，后面的子图检索和多跳排序才不会一直在碎图上打补丁。

\subsection{GraphRAG 落地时最难的不是建图，而是子图召回与排序}
旧稿在 GraphRAG 章节里还有一个很有价值的工程提醒：真正让图检索难落地的，常常不是“如何把实体和关系存起来”，而是\textbf{当用户提问时，到底该捞出哪一小块子图、再按什么顺序交给模型。}

这一步至少包含三层判断：
\begin{itemize}
  \item \textbf{实体抽取够不够稳}：如果“报销单”“差旅申请”“费用申请单”抽成了不同实体，图本身就会碎裂。
  \item \textbf{子图半径该多大}：半径太小，多跳关系不够；半径太大，噪声和成本迅速上升。
  \item \textbf{图证据怎么排序}：即使捞对了子图，也不代表所有节点和边都值得同等进入上下文。
\end{itemize}

因此，GraphRAG 工程上通常也需要粗排与精排。先用较轻规则或图统计特征筛掉大部分无关节点，再用更贵的语言模型或重排序器判断哪些关系真正支撑当前问题。这样讲，GraphRAG 就不再是一个抽象的“有图就更强”，而是一套比普通 RAG 更重的结构检索预算分配机制。

\begin{figure}[H]
\centering
\begin{tikzpicture}[
  font=\small,
  box/.style={llm box, minimum width=3cm, minimum height=1cm},
  arr/.style={llm arr}
]
\node[box] (parse) {文档理解\\PDF / OCR / 表格};
\node[box, right=0.9cm of parse] (split) {结构保留\\切分与元数据};
\node[box, right=0.9cm of split] (query) {查询增强\\改写 / HyDE / 多查询};
\node[box, right=0.9cm of query] (rerank) {召回优化\\重排 / 负样本};
\node[box, below=1.1cm of query] (graph) {GraphRAG\\强关系与多跳问题};
\draw[arr] (parse) -- (split);
\draw[arr] (split) -- (query);
\draw[arr] (query) -- (rerank);
\draw[arr] (query) -- (graph);
\end{tikzpicture}
\caption{文档理解与召回增强的推荐优化顺序}
\label{fig:rag-optimization-stack}
\end{figure}

\subsection{多向量与多模态 RAG：图片、版面和表格都可以成为证据}
旧稿里还有一部分内容提醒我们：企业知识并不总是纯文本。很多关键信息藏在截图、流程图、拓扑图、表单样例、版面布局甚至带箭头的操作说明里。若系统只把纯文本切块入库，这些信息就会天然缺席，最后看起来像“模型没学会”，其实是证据从来没进过知识库。

一个更实用的做法，是把同一份知识表示成多种可检索视角。例如，同一页 PDF 可以同时保留标题路径、正文块、表格结构、图片说明和版面区域摘要；同一张图片既可以有 OCR 文本，也可以有人工或模型生成的图像摘要。这样做的意义，不在于把所有问题都升级成多模态，而在于让“图片里才有的关键信息”至少有机会被召回。

多向量检索器在这里很有用，因为它允许一个知识对象不只对应一个 embedding。对企业知识助手来说，一条制度可以同时有“标题向量”“摘要向量”“FAQ 问句向量”；一张表可以有“表头向量”和“关键行向量”；一张截图可以有“OCR 文本向量”和“图像说明向量”。真正回答时，系统再把命中的表示映射回原始文本块、表格块或图像区域，而不是直接拿抽象表示去生成答案。

\subsection{自适应 RAG：把检索做成条件动作，而不是固定流水线}
旧稿里还有一类主题在新版里偏薄，那就是 Self-RAG、模块化 RAG 这类“动态检索控制”方法。它们共同强调的一点是：\textbf{不是每个问题都该走完全一样的检索链。} 有的问题模型本来就能回答，有的问题只需一次检索，有的问题则需要“先检索一轮，再批判证据是否够用，必要时再检索第二轮”。

从工程角度看，自适应 RAG 主要解决三种情况：
\begin{itemize}
  \item \textbf{并不是所有请求都值得付同样检索成本}：高频 FAQ 或明显短问题，可能根本不需要走最重的检索链；
  \item \textbf{检索到不等于证据够用}：第一轮召回可能相关，但不一定足以支撑最终回答；
  \item \textbf{系统需要显式决定“何时停手”}：否则多轮检索会无限堆成本，却未必持续带来收益。
\end{itemize}

Self-RAG 可以理解为把“是否检索、如何批判证据、是否继续检索”也纳入模型决策；模块化 RAG 则更像把问题改写、检索、重排序、证据批判和回答生成拆成可替换的部件。两者都不是最小闭环阶段的第一优先级，但当系统已经跑通静态检索链，却仍然频繁出现“证据不够却硬答”“简单问题也走重链路”时，它们就会开始有意义。

落到工程里，一条够用的动态链路通常长这样：先做问题分流，判断它属于简单 FAQ、普通检索题还是复杂综合题；再决定走轻检索还是重检索；拿到候选后先做一次证据批判，检查当前材料是否足以支撑回答；若证据仍不足，再决定是否触发第二轮检索、查询扩展或直接拒答。这样一来，自适应 RAG 就不再只是一个抽象名词，而变成了“把检索预算花在该花的地方”的工程机制。

\subsection{备用模型与安全兜底：不是所有问题都交给同一条链}
旧稿的系统化痛点章节还提醒了一类很现实的工程需求：\textbf{当主链证据不够、解析失败、结构化输出不稳或问题风险升高时，系统要知道谁来接手。} 这并不一定意味着再上一个更大的模型，更多时候意味着准备好一条兜底路径。

常见的兜底方式包括：
\begin{itemize}
  \item 文档解析失败时，回退到更保守的 OCR 或人工标注版本；
  \item 主检索链证据不足时，切到更重的多查询或人工审核链；
  \item 高风险问题触发规则校验、权限确认或明确升级人工；
  \item 结构化输出失败时，先保留证据与结论，再由校验器或次级模型补格式。
\end{itemize}

这部分之所以重要，是因为一套真正上线的 RAG 系统，评估的不只是“主链在理想条件下能不能答对”，还包括“输入脏了、证据缺了、问题变难了以后，它会不会有体面的失败方式”。对企业知识助手来说，能安全地说“证据不足，建议升级处理”，往往比硬把所有问题都塞进同一条主链更成熟。

\subsection{chunk size、overlap 与 top-k 要靠验证集调，不要靠直觉}
旧稿里有大量关于 \texttt{chunk\_size}、\texttt{chunk\_overlap}、\texttt{top-k} 的实验片段，这些东西不能只留下参数名而没有方法论。真正稳定的 RAG 调参，不是“觉得 512 很常见就先上 512”，而是先准备一套小而硬的验证集，再用这套验证集去比较不同组合。

一套够用的调参协议通常包括：
\begin{enumerate}
  \item 固定一组代表性问题，至少覆盖短问句、长问题、表格问答、跨段落综合题和拒答题。
  \item 对 \texttt{chunk\_size}、\texttt{overlap}、\texttt{top-k}、重排序阈值等参数做小范围网格搜索，而不是一次改很多层策略。
  \item 同时记录召回质量、答案忠实度、上下文长度、首字时间和单位请求成本，而不是只看“答对率”。
  \item 优先选择\textbf{能满足质量下限的最小上下文方案}，而不是无上限把更多片段塞进 prompt。
\end{enumerate}

如果团队一开始没有足够人工问题，也可以先用日志问题、FAQ、制度标题改写和少量合成问题搭一个启动集，但一定要在关键参数收敛后补人工抽查。因为合成问题通常比真实问题更“配合”文档表达，很容易高估检索效果。

教材里最该保留的不是某个参数最佳值，而是这条纪律：\textbf{RAG 参数必须和评测集版本一起被记录。} 否则今天把 \texttt{top-k} 从 3 调到 8 看起来变好了，明天文档切分方式一改，所有结论都可能失效，却没人说得清到底是参数问题、切分问题还是评测样本变了。

\subsection{先查覆盖率，再查召回率：知识库里没有的答案别逼系统编}
旧稿的痛点分析里还有一条很应该留下的纪律：不是所有答错都值得继续调检索。很多时候，系统回答不稳并不是因为检索器太弱，而是因为\textbf{知识库里压根没有当前问题需要的那份权威材料，或者材料已经过期。}

这在企业环境里尤其常见。制度刚更新但新版本还没入库，区域细则存在于附件却没被抽取，权限文档在私有空间里根本没同步，或者同一规则同时存在“正式制度版”和“群里流传截图版”。如果这些覆盖问题没有先查清楚，团队就很容易把“没有答案”误诊成“没召回到答案”，然后在错误方向上持续加复杂度。

一个更成熟的排查顺序通常是：
\begin{enumerate}
  \item 先确认当前问题所需的权威文档是否真的在知识库里；
  \item 再确认最新版本、附件页、图片页和表格页是否都被正确入库；
  \item 再确认是否存在权限隔离、地域隔离或业务线隔离导致的候选缺失；
  \item 只有在这些都成立以后，才继续查召回、排序和生成链。
\end{enumerate}

这条顺序之所以重要，是因为检索优化永远无法创造不存在的证据。对企业知识助手来说，能及早发现“当前知识库覆盖不到这个问题”，并把它回流到数据治理或知识运维流程里，本身就是系统成熟度的一部分。

\subsection{先从症状判断该修哪里}
为了避免优化路径失控，我们可以反过来从症状判断入口。如果召回片段本身就乱，优先查解析。如果召回片段很像问题却答不上来，优先查重排序和难负例。如果问题一换说法就命中率大幅下降，优先查查询增强。如果问题天然依赖组织关系或依赖链，再考虑图结构。

\begin{table}[H]
\centering
\small
\begin{tabularx}{\textwidth}{>{\raggedright\arraybackslash}p{3.1cm}>{\raggedright\arraybackslash}p{3.2cm}>{\raggedright\arraybackslash}X}
\toprule
\textbf{观察到的现象} & \textbf{优先怀疑哪一层} & \textbf{第一动作} \\
\midrule
片段混入页眉页脚、断裂句子、乱码表格 & 文档理解层 & 先修 PDF 解析、OCR 和版面顺序，不急着换模型 \\
召回“像相关”但不能支持答案 & 排序层 & 先上重排序，检查是否被难负例误导 \\
用户换一种说法就检索不到 & 查询层 & 先试问题改写、HyDE 或多查询扩展 \\
问题需要跨实体、跨步骤找依赖链 & 结构层 & 再考虑 GraphRAG 或关系抽取 \\
\bottomrule
\end{tabularx}
\caption{从症状倒推 RAG 优化入口}
\label{tab:rag-symptom-routing}
\end{table}

\section{主案例推进}
在企业知识助手的主案例里，最小闭环跑通后，团队很快发现两个新的现实问题。第一，制度 PDF 中的表格经常被切碎，导致“住宿标准”和“超标审批”无法一起被解释。第二，员工的提问非常口语化，常用说法与制度标题并不一致，导致召回结果不稳定。

\subsection{第一轮优化先修输入，不先修模型}
我们先抽样看错误日志，发现很多失败案例里，召回片段本身就不完整：页眉页脚混入正文，双栏 PDF 顺序错位，表格只有数据没有表头。这个现象说明，当前瓶颈根本不在模型，而在输入理解。

于是团队先做三件事：
\begin{enumerate}
  \item 对制度类 PDF 保留标题层级和页码；
  \item 对表格单独抽取结构块，而不是和普通段落混切；
  \item 对双栏文档重排阅读顺序，避免左右栏交叉串行。
\end{enumerate}

这一轮调整之后，系统不一定立刻“回答更聪明”，但它开始“看见更正确的材料”。对 RAG 来说，这是更基础也更关键的进步。

\subsection{第二轮优化再修用户语言与文档语言的错位}
输入修干净后，另一类失败仍然存在：员工问“住宿超标还能报吗”，系统能召回“住宿标准”，却漏掉“超标审批例外”。这说明文档里有答案，但用户表达和制度表达没有完全对上。

此时我们引入轻量问题改写，并对高频口语问题增加 HyDE 检索路径。前者把“还能报吗”改写成“住宿超标审批与报销规则”，后者则通过假想答案把“住宿标准”“审批例外”“附件要求”等相关术语一起带入检索空间。由于这里只是在修表达落差，所以投入和收益都比较合适。

\subsection{第三轮优化才考虑排序和难负例}
再往后，系统出现另一种现象：召回列表里经常已经有正确片段，但前排也混入了很多“差旅相关但不是住宿审批”的材料，例如交通补助说明、出差借款流程、审批附件要求。此时继续增大 \texttt{top-k} 只会让噪声更多，于是团队开始引入重排序。

如果重排序后仍然有大量“相关但无用”的片段排在前面，下一步才会考虑用领域样本去挖掘难负例，例如把“住宿标准”和“餐补标准”“权限审批”和“权限备案”这类高相似但语义不同的片段成对加入训练或评估集。

\begin{table}[H]
\centering
\small
\begin{tabularx}{\textwidth}{>{\raggedright\arraybackslash}p{2.6cm}>{\raggedright\arraybackslash}p{3cm}>{\raggedright\arraybackslash}p{2.9cm}>{\raggedright\arraybackslash}X}
\toprule
\textbf{失败症状} & \textbf{采取动作} & \textbf{为什么先做它} & \textbf{停手条件} \\
\midrule
PDF 片段顺序混乱 & 修版面解析与结构切分 & 输入脏时，后续所有检索都在坏文本上进行 & 片段已基本可读、可回到原页 \\
表格信息经常缺字段 & 单独保留表头与行列关系 & 表格语义首先取决于结构绑定 & 表格问答不再频繁漏列名和限制条件 \\
口语问题命中率低 & 问题改写或 HyDE & 用户语言与文档语言存在明显落差 & 同类口语问法召回趋于稳定 \\
结果“相关但不够对” & 重排序或难负例治理 & 候选已出现正确文档，但排序质量不足 & 关键证据能稳定进入前几位 \\
\bottomrule
\end{tabularx}
\caption{企业知识助手在 RAG 第二阶段的优化动作与停手条件}
\label{tab:rag-optimization-actions}
\end{table}

\subsection{不是所有关系问题都要立刻上 GraphRAG}
项目推进到这里时，团队内部开始提出：既然企业内部有组织架构、系统依赖和审批链，是不是应该马上上 GraphRAG？这个问题不能只靠“技术先进不先进”来判断，而要看主问题是不是已经表现出明显的多跳关系需求。

在我们的主案例里，大部分高频问题仍然是制度问答和说明文档问答，本质是文档查证，不是图上推理。因此，GraphRAG 暂时只作为局部能力储备，而不是全局底座。等到后续确实出现“某权限申请依赖哪几个系统负责人逐级审批”“某服务故障影响哪些下游业务”这类强关系问题时，再引入图结构会更合理。

\subsection{第四轮优化：只让复杂问题走重检索链}
当静态检索链逐步稳定后，团队又遇到一个新现象：很多简单 FAQ 问题也被送进了完整的“改写 + HyDE + 多召回 + 重排序”链路，导致成本偏高；而少数复杂问题虽然进了检索链，却仍然因为第一轮证据不足而答得很勉强。

于是第四轮优化不再只是“再加一种检索技巧”，而是引入更轻量的自适应控制：
\begin{enumerate}
  \item 对高频、短问题先做是否需要重检索的判别；
  \item 对低置信度回答增加证据批判步骤，判断当前片段是否真能支持答案；
  \item 只有在证据不足时，才触发第二轮检索或更重的查询增强。
\end{enumerate}

这一步没有让系统“突然更聪明”，但它让资源分配更合理了：简单问题不再平白付出重检索代价，复杂问题也终于有机会在证据不足时补一次检索，而不是直接把第一轮半对不对的结果包装成完整答案。

\subsection{第五轮优化：把版本过滤和混合检索接到主链上}
系统上线一段时间后，团队又遇到一个很典型的新问题：同样在问“差旅住宿上限”，上海员工和海外员工拿到的答案不一样；同样在问“代码库权限”，有的人需要的是旧流程，有的人需要的是新审批单。排查后发现，问题不在模型突然变差，而在于知识库里同时存在多版本、多地域、多业务线材料，而主链还没有把这些结构化约束前移。

于是第五轮优化做了两件事。第一，在检索入口增加地域、制度生效日期、员工类别和文档版本过滤，让候选先落进正确范围；第二，对包含制度名、错误码和产品名的问题切到混合检索，让关键词路径负责精确命中，向量路径负责容纳口语表达。这样一来，系统不必再让大模型在一堆“都像相关”的材料里猜哪一份更适用。

这一轮调整有一个很重要的启发：当问题里已经带着明确业务约束时，最值钱的优化往往不是再上一层更重的生成技巧，而是\textbf{把候选池先收对。} 对企业知识助手来说，先把“哪份文档有资格参与竞争”判断清楚，往往比让模型在错误候选中做更聪明的选择更有效。

\section{常见坑与工程取舍}
\begin{itemize}
  \item \textbf{解析错误却去调 embedding}：这是 RAG 项目里最常见的误诊之一。输入坏了，换模型通常只是换一种方式读错。
  \item \textbf{把所有查询增强都默认打开}：多查询和 RAG-Fusion 可能带来收益，但也会带来延迟、成本和噪声，不应无差别启用。
  \item \textbf{把 GraphRAG 当成通用升级包}：如果知识本身没有明显关系结构，引入图只会增加构建和维护复杂度。
  \item \textbf{只看召回率，不看证据质量}：召回更多不等于召回更好，企业问答更怕“相关但错误指向”的噪声片段。
  \item \textbf{一次同时改解析、切分、检索和提示}：这样做即使效果变好，也很难知道到底哪一项真正起作用。
  \item \textbf{把 overlap 开得很大却不做去重}：这会让重复候选淹没真正的新证据，看起来“更全”，实际上更乱。
  \item \textbf{把版本号、生效日期和权限范围留给模型自己猜}：这些本该在检索入口通过元数据过滤解决，而不是放到生成阶段碰运气。
\end{itemize}

再补四个在企业文档里特别典型的取舍：
\begin{itemize}
  \item \textbf{优先保留结构还是优先追求切分简单}：在制度、表格和流程文档里，通常优先保留结构。
  \item \textbf{优先上查询增强还是优先上重排序}：如果同义问法一多就失灵，先查查询增强；如果正确片段已常常出现但排不前，先查重排序。
  \item \textbf{优先做图结构还是继续做文档治理}：只要高频问题仍是文档查证，文档治理通常比图结构更值得先投。
  \item \textbf{优先提覆盖还是优先提稳定}：进入优化期后，稳定命中高频问题通常比扩更多边缘问题更有价值。
\end{itemize}

\section{本章练习}
\begin{exercise}[文档解析失败]
某制度 PDF 的召回片段中经常混入页眉页脚和断裂表格，导致答案看起来“引用了材料”但实际不可读。此时最优先的优化方向是什么？为什么不该先去调查询增强？
\end{exercise}

\begin{solution}
最优先的方向是{\heiti 文档理解与结构保留}，也就是先修 PDF 解析、OCR、阅读顺序和表格抽取。因为当前问题发生在检索之前，系统连“正确文本是什么”都还没稳定拿到。查询增强解决的是用户表达和文档表达的落差，无法补救已经被解析坏的输入。
\end{solution}

\begin{exercise}[何时考虑 GraphRAG]
什么样的问题更适合考虑 GraphRAG，而不是继续叠加问题改写、HyDE 或 RAG-Fusion？
\end{exercise}

\begin{solution}
当问题天然依赖{\heiti 实体关系和多跳链路}时，更适合考虑 GraphRAG。例如审批链依赖、组织结构依赖、系统组件依赖、告警传播链和供应关系等。这类问题的关键不是“文档里有哪句话”，而是“谁和谁有什么关系、哪一步依赖哪一步”。如果问题仍然主要是在制度文本里查证条款，通常先把文档理解、结构保留和召回质量做好更划算。
\end{solution}

\section{延伸阅读}
\begin{itemize}
  \item \textbf{HyDE、RAG-Fusion 相关论文与开源实现}：适合理解查询增强为什么能弥合用户语言和文档语言之间的落差。
  \item \textbf{Unstructured、MinerU、Docling、Nougat 等文档解析工具资料}：适合理解 PDF、扫描件和表格解析的工程现实。
  \item \textbf{重排序与检索器微调资料}：适合理解“相关但无用”的召回为什么需要专门治理。
  \item \textbf{GraphRAG 官方材料与案例}：适合理解图结构真正适合解决什么类型的问题，而不是把它当作通用组件。
\end{itemize}

\section{小结}
RAG 优化不是一串彼此独立的小技巧，而是一条有先后顺序的输入理解与召回增强路线。真正有工程价值的优化，必须能回答三个问题：它修的是哪一层、为什么现在修这层、什么情况下可以先停手。只要把这个顺序守住，企业知识助手就不会滑向“技术名词堆叠”，而会稳定地从最小闭环迈向可运营系统。

\section{下一章过渡}
当文档理解和召回质量逐步稳定后，团队会遇到新的问题：到底是模型层、系统层还是应用层在变好？下一章将建立评测体系，把模型效果、RAG 质量、系统稳定性和人工反馈统一纳入度量，让后续优化不再只靠感觉推进。
